\section{Discussion}

\pending{Opening summary of answers to RQ1--RQ4.}

\subsection{Planning horizon and re-optimization frequency}
\pending{Interpretation against rolling-horizon theory \citep{baker1979analytic,sethi1991theory,
sahin2008master} and forest operations studies \citep{beaudoin2008hierarchical,frisk2016detailed,
arora2023detailed}. State explicitly which conclusions of \citet{jaffray2026thesis} hold and which
are revised: (a) re-optimization frequency dominates solution quality; (b) horizons beyond about
8 weeks give diminishing returns; (c) weekly re-optimization over 4--8-week horizons is a practical
sweet spot; (d) horizon length mainly drives computing effort.}

\subsection{Why state carry-forward and stitched-plan evaluation matter}
\pending{Rolling-horizon evaluations must carry the operational state (remaining volume, staged
inventory between phases, machine positions) between windows and score the implemented stitched
schedule. Evidence: the size and direction of the differences from the earlier evaluation, without
reporting its values as results.}

\subsection{Solver choice and tractability}
\pending{Cold-start exact solves on large windows; empty incumbents; the earliness stage and work
deferred past the locked span; the case for heuristics or warm-started MIP at operational scale
\citep{frisk2016detailed,audy2023planning}.}

\subsection{Practical implications}
\pending{Guidance on configuring rolling-horizon planning, carried over from the thesis
discussion where the new results support it: directional long-range plans
\citep{myers2004analysis,beaudoin2008hierarchical}; open tools lowering the perceived risk of
changing practice \citep{de2023perceptions}; dependence on scale and context.}

\subsection{Limitations}
This evaluation is deterministic. It isolates the structural effects of the rolling-horizon
configuration and leaves out weather, breakdowns, and forecast error, which are the usual reason
for re-planning. Extending it with the stochastic playback in FHOPS is the natural next step. The
case studies are semi-synthetic. They use real BC inventory and terrain data but one harvesting
system, synthetic fleets, one landing per block, no mobilisation costs, and no block time windows.
SA was run with one seed and a fixed iteration budget, so its results show the behaviour of one
trajectory rather than a distribution. Time-limited MIP solves are not proven optimal, and their
results depend on the solver, the limit, and the hardware. FHOPS windows place no value on work in
progress at the end of a window. A time-limited earliness stage can therefore leave work deferred
beyond the locked span, which penalises long windows with short lock spans. FHOPS~1.0.1 can also
export assignments with negligible planned production (below $10^{-3}$~m$^3$), which replay as
sequencing violations of negligible volume \pending{reference FHOPS issue and the
threshold used}. Absolute computing times depend on hardware. \pending{Add or remove limitations
once the results are known.}

\section{Conclusions}
\pending{Three to four sentences answering the research questions, plus future work: uncertainty,
real operational data, other harvesting systems and regions.}
